% paper/sections/advisor.tex — Grounded AI Advisor and Scientific Provenance
% Written in Phase 88 (Plan 04). TikZ advisor-flow diagram added in Phase 90.
% Key entry points: build_diagnostics, advise, auto_tune, fdars_method_references MCP tool

\section{Grounded AI Advisor and Scientific Provenance}
\label{sec:advisor}

A recurring challenge in applied functional data analysis is \emph{parameter
selection}: bandwidth for kernel smoothing, the number of basis functions, the
number of FPCA components, or the number of clusters.  While numerical
criteria (GCV, AIC, silhouette score) narrow the range, interpreting their
output in the context of a specific analytical goal --- and translating that
interpretation into a revised parameter setting --- requires domain and
methodological knowledge.  \texttt{fdars} addresses this through a
\emph{grounded parameter advisor} paired with a \emph{scientific-provenance
layer}, both absent from peer FDA packages (see Table~\ref{tab:comparison} and
the evidence in \texttt{paper/comparison\_evidence.md}).

\begin{figure}[ht]
\centering
\resizebox{\linewidth}{!}{%
\begin{tikzpicture}[
  font=\small,
  box/.style={draw, rounded corners=3pt, minimum width=2.8cm, minimum height=0.95cm,
              align=center, inner sep=5pt},
  prov/.style={draw, rounded corners=3pt, minimum width=3.8cm, minimum height=0.95cm,
               align=center, inner sep=5pt, fill=purple!10, dashed},
  arr/.style={-{Stealth[length=5pt]}, thick},
  node distance=1.0cm and 1.6cm,
]
  \node[box, fill=gray!15] (result) {\texttt{fdars}\\result dict};
  \node[box, fill=blue!10, right=of result] (diag) {\texttt{build\_diagnostics}};
  \node[box, fill=orange!10, right=of diag] (advise) {\texttt{advise}};
  \node[box, fill=orange!15, below=1.4cm of advise] (tune) {\texttt{auto\_tune}};
  \node[prov, below=1.9cm of diag] (prov)
    {provenance map\\\texttt{fdars\_method\_references}};

  \draw[arr] (result) -- (diag)
    node[midway, above, font=\footnotesize] {result};
  \draw[arr] (diag) -- (advise)
    node[midway, above, font=\footnotesize] {diagnostics};
  \draw[arr, dashed] (advise) -- (tune)
    node[midway, right, font=\footnotesize] {loop};
  \draw[arr, dashed] (tune) to[out=200,in=-20] (diag);
  \node[box, fill=gray!10, left=of prov] (mcp) {MCP client\\(query by callable)};
  \draw[arr, dashed] (prov) -- (mcp)
    node[midway, above, font=\footnotesize] {refs};

  \node[font=\footnotesize\itshape, above=4pt of diag] {LLM-free};
  \node[font=\footnotesize\itshape, above=4pt of advise] {provider-agnostic LLM};
\end{tikzpicture}%
}
\caption{Advisor and scientific-provenance data flow.
\texttt{build\_diagnostics} is entirely deterministic (no LLM);
\texttt{advise} queries a user-configured language-model provider using only
the structured diagnostics as context;
\texttt{auto\_tune} iterates both in a bounded closed-loop.
The provenance map (\texttt{\_references\_map.json}) feeds
\texttt{fdars\_method\_references}, an MCP tool that resolves callable names
to peer-reviewed papers at query time; this path is independent of the
advisor loop.}
\label{fig:advisor}
\end{figure}

\subsection{The Grounded Advisor}

The advisor's three core entry points in \texttt{fdars.advisor} are listed
below; further helpers include \texttt{compare\_methods} (comparative method
selection over several candidate diagnostics) and \texttt{pipeline\_report}:

\sigheading{\texttt{build\_diagnostics} --- the LLM-free grounding step.}
\begin{lstlisting}[style=fdarsinput]
build_diagnostics(
    result, method, *, argvals=None, ...
)
\end{lstlisting}
    It accepts the result dictionary
    returned by an \texttt{fdars} analysis function (e.g.\ the dict returned
    by \texttt{Fdata.to\_pc}) and a method identifier string (e.g.\
    \texttt{"fpca"}, \texttt{"smoothing"}, \texttt{"clustering"}) and
    returns a structured diagnostics dictionary derived entirely from
    deterministic computation on the result.  The content of the diagnostics
    dict is method-specific: for \texttt{fpca} it includes the eigenvalues,
    the component count, per-component and cumulative explained-variance
    ratios expressed as shares of the total variance (computed from the
    result's centred data and quadrature weights; a
    \texttt{variance\_denominator} field records when a result lacks these
    and the ratios fall back to shares of the retained components), and a
    phase-leakage indicator (the share of retained variance outside the first
    component); for \texttt{smoothing} it includes the GCV curve and its
    optimum, effective degrees of freedom, and AIC/BIC approximations; for
    \texttt{clustering} it includes $k$, cluster sizes and means, and
    amplitude and phase separation between clusters.  No
    language model is involved; the diagnostics are fully reproducible given
    the same result input, and can be inspected, logged, or tested
    independently of the LLM step.  Supported method identifiers include
    \texttt{alignment}, \texttt{basis}, \texttt{classification},
    \texttt{clustering}, \texttt{depth}, \texttt{fpca}, \texttt{frechet},
    \texttt{fts}, \texttt{inference}, \texttt{outliers}, \texttt{regression},
    \texttt{regression\_cv}, \texttt{represent}, \texttt{scoring},
    \texttt{smoothing}, and \texttt{spm}.  The capability tour
    (Section~\ref{sec:capabilities}) demonstrates \texttt{build\_diagnostics}
    on an FPCA result.

\sigheading{\texttt{advise} --- the provider-agnostic LLM entry point.}
\begin{lstlisting}[style=fdarsinput]
advise(
    diagnostics, *, task, domain_context,
    provider=None, model=...,
) -> Advice
\end{lstlisting}
    It accepts the diagnostics
    dict produced by \texttt{build\_diagnostics} --- never raw data or a
    result dict directly --- and calls a user-configured language-model
    provider (Anthropic, OpenAI or an OpenAI-compatible endpoint, Google
    Gemini, or a local Ollama server) to generate parameter recommendations
    framed in terms of the supplied task description and domain context.  The
    provider defaults to Anthropic unless overridden by the
    \texttt{provider} argument or the \texttt{FDARS\_ADVISOR\_PROVIDER}
    environment variable; the model is taken from the \texttt{model}
    argument, else the \texttt{FDARS\_ADVISOR\_MODEL} environment variable,
    else the chosen provider's default (\texttt{claude-opus-4-8} for
    Anthropic).  API keys are read
    from the user's environment and are never embedded in the library.  Because the LLM receives only the structured,
    deterministic diagnostics as its context, its recommendations trace back
    to computed facts rather than free-form interpretation of raw arrays.
    This \emph{grounding invariant} is what distinguishes the advisor from a
    plain LLM chat interface applied to a data summary: the diagnostics
    constitute an explicit, auditable evidence layer that the language model
    reasons over, and the same diagnostics can be separately inspected by the
    user to verify that the LLM's advice is consistent with the numbers.
    The invariant is also enforced mechanically: before \texttt{advise}
    returns, a grounding check extracts every numeric value cited in the
    advice's evidence and requires it to equal a diagnostic value at the
    citation's own precision; otherwise the call fails, with no retry.
    Qualitative statements pass unchecked, so the guard stops fabricated
    statistics but cannot verify the reasoning that connects them.

\sigheading{\texttt{auto\_tune} --- closed-loop parameter search.}
\begin{lstlisting}[style=fdarsinput]
auto_tune(
    dataset_id, method, *, target_metric=None,
    max_steps=10, guard=True, ...
)
\end{lstlisting}
    A closed-loop parameter search that iterates
    \texttt{build\_diagnostics} $\to$ \texttt{advise} $\to$ re-fit in a
    bounded loop (at most \texttt{max\_steps} iterations).  With the
    default \texttt{guard=True}, a method-specific guard re-checks one
    secondary diagnostic after every re-fit---FPCA tuning stops if the
    phase-leakage indicator exceeds 0.5, smoothing tuning if the effective
    degrees of freedom exceed $0.9\,n$, basis-penalty tuning if the optimal
    GCV score degrades by more than 20\%, and clustering tuning on a cluster
    of size below 2---so the loop cannot improve its target metric by
    violating that guarded diagnostic; other quality measures are not
    checked.  The loop stops when the
    step budget is exhausted, when the guard is violated, when proposed
    parameters revisit or oscillate, or after a run of non-improving steps;
    every number it acts on is computed by \texttt{fdars}.

The design separates the deterministic grounding step (\texttt{build\_diagnostics})
from the probabilistic recommendation step (\texttt{advise}), ensuring that
the pipeline can be tested and audited at the boundary between the two.
In particular, \texttt{build\_diagnostics} can be unit-tested against known
result dicts without any LLM dependency, and the diagnostics dict format is
stable across provider changes.

\subsection{How the Advisor Is Evaluated---and What Is Not Evaluated}

The advisor's deterministic parts are tested offline in CI with a stub
language-model provider.  An offline evaluation suite checks, on synthetic
diagnostics with a known best value, that comparative method selection
(\texttt{compare\_methods}) returns the best-scoring method, and, with
injected diagnostics that improve by construction, that the tuning loop
accepts improving steps and terminates within its budget; neither test fits
real data.  A
separate regression suite exercises the grounding check on valid citations
(rounded values, negative numbers, indices, scientific notation) and on
fabricated ones, which must be rejected.  A live-provider smoke test exists
but runs only when an API key is supplied, so it is not part of CI.

We have \emph{not} measured the quality or usefulness of the recommendations
a real language model produces---there is no benchmark against expert choices
and no user study---and the advice may vary between providers and model
versions.  The contribution claimed here is therefore the grounded
architecture and its enforced invariant, not demonstrated recommendation
quality.

\subsection{Scientific-Provenance Layer}

\texttt{fdars} ships a machine-readable provenance map,
\texttt{python/fdars/\_references\_map.json}, that links public callables
(243 of the \ncallables{} at present) to candidate foundational papers from
which their algorithms derive.  Besides 10 placeholder records grouping
not-yet-attributed callables, the map contains \ndocpapers{} candidate papers
(44 with DOIs or arXiv identifiers), of which 6 are author-verified; at
present \ncoverage{} of the \ncallables{} callables (\npubliccallables{} module
callables plus the 28 \texttt{Fdata} entries: \nfdata{} public methods and the
constructor) have at least one author-verified paper entry, so coverage is
deliberately partial and honest rather than complete.
Each paper record carries the title, authors, year, DOI and/or URL,
publication type, the \texttt{fdars} callables that implement it, notes, and
a \texttt{curated} boolean indicating whether the entry has been manually
verified against the primary source; where available (24 of the \ndocpapers{} candidate papers)
it also carries cross-language implementation pointers (R, Python,
occasionally MATLAB) with confidence ratings.  Entries marked
\texttt{curated: false} are present as cross-reference candidates but are
explicitly flagged so that downstream tooling can distinguish author-verified
citations from unverified suggestions.

The provenance map is served at runtime through the
\texttt{fdars\_method\_references} tool of the package's Model Context
Protocol (MCP) server \citep{anthropic_mcp_2024}, \texttt{fdars.mcp}.
Given a callable name, the tool returns the associated paper entries, DOIs, and
cross-language pointers --- making the provenance queryable by any
MCP-compatible development environment (e.g.\ Claude Desktop, VS Code Copilot)
without requiring a network call at query time, since the full map is bundled
with the package.  The same map is used to generate the References
documentation page of the MkDocs site (\texttt{docs/references.md}), so the
paper-level provenance visible to readers and the provenance served through
the MCP server are built from one source of truth (regenerating the page is
a manual step).

\subsection{Differentiator}

The evidence dossiers in \texttt{paper/comparison\_evidence.md} confirm that
none of the peer FDA packages surveyed --- scikit-fda, FDApy, R \texttt{fda},
\texttt{fda.usc}, \texttt{refund}, \texttt{funData}/\texttt{tidyfun}, or
Matlab \texttt{fdaM}/PACE --- provides either a grounded parameter advisor or
a machine-readable scientific-provenance layer (see the ``Grounded advisor +
scientific provenance'' row in Table~\ref{tab:comparison}, which shows
\checkmark{} only for \texttt{fdars}).  These two contributions address
distinct barriers to adoption: the advisor lowers the methodological cost of
parameter selection for applied users, while the provenance layer gives
researchers a traceable link from any \texttt{fdars} result back to the
peer-reviewed literature from which the method derives.
